% 4x4、8级灰度的教学直方图：均衡前计数(2:5,3:6,4:3,5:2)，按 s_k=round(7*cdf) 映射为 {2,5,6,7}。
% 两个面板并排，总宽约 92mm。
\begin{tikzpicture}[x=4.2mm,y=2.6mm,font=\figurefont\small]
  \foreach \ox/\title/\data in {0/均衡前/{2/5,3/6,4/3,5/2},11.5/均衡后/{2/5,5/6,6/3,7/2}}{
    \begin{scope}[xshift=\ox*4.2mm]
      \foreach \y in {2,4,6}{
        \draw[FigureGrid,line width=.4pt] (-.6,\y)--(7.6,\y);
        \node[anchor=east,font=\figurefont\footnotesize] at (-.7,\y) {\y};
      }
      \foreach \k/\n in \data{
        \fill[FigureBlueFill,draw=FigureBlue,line width=.6pt] (\k-.38,0) rectangle (\k+.38,\n);
        \node[above,font=\figurefont\footnotesize] at (\k,\n) {\n};
      }
      \draw[FigureStroke,line width=.6pt] (-.6,0)--(7.6,0);
      \foreach \k in {0,...,7}{\node[below,font=\figurefont\footnotesize] at (\k,0) {\k};}
      \node at (3.5,-2.6) {灰度级};
      \node at (3.5,8.3) {\title};
    \end{scope}
  }
  \draw[-{Stealth[length=2mm]},FigureStroke,line width=.8pt] (8.1,3.5)--(10.2,3.5);
  \node[above,font=\figurefont\footnotesize,text=FigureMuted] at (9.15,3.5) {$T$};
\end{tikzpicture}
